\section{The compact family and its rank-four bundle}\label{sec:geometry}

We first construct the degeneration and its bundle on the entire compact
family. The smoothing parameter and the unitary gluing parameter will vary
independently. Both are fixed during each metric-flow trajectory.

\subsection{An elliptic family with four central components}
Let $\Delta$ be a disc about zero of radius less than one. In
$\PP^2\times\Delta$ consider the family
\begin{equation}\label{eq:cubic}
 \mathcal W=\{y^2z=x(x-z)(x-t^2z)\}.
\end{equation}
Its section $[0:1:0]$ supplies the origin on every smooth elliptic fibre.
The involution is $\sigma(x,y,z,t)=(x,-y,z,t)$.

\begin{proposition}\label{prop:elliptic-family}
Two successive blowups at singular points over $t=0$ give a proper
projective flat family $\cE\to\Delta$ with smooth total surface. Its
noncentral fibres are smooth elliptic curves. Its central fibre is a cycle
$D_0\cup D_1\cup D_2\cup D_3$ of smooth rational curves, with local
equation $uv=t$ at each of its four nodes. The involution lifts to $\cE$,
acts by $D_j\mapsto D_{-j}$, and fixes no node. Indices are taken modulo four.
\end{proposition}

\begin{proof}
For $t\ne0$ the roots $0,1,t^2$ are distinct, so the fibre is a smooth
plane cubic. In the chart $y=1$, the derivative with respect to $z$ of the
defining equation is one at $x=z=0$, proving smoothness at the section.
In the chart $z=1$, write
$f=y^2-x(x-1)(x-t^2)$. The identities
$f_y=2y$ and $f_t=2tx(x-1)$ show that the only singular point of the total
surface over the chosen disc is $(x,y,t)=(0,0,0)$. At $t=0$ this also
follows directly from $f=f_x=f_y=0$.

Choose the square root of $1-x$ equal to one at zero and set
\begin{equation}\label{eq:a3-coordinates}
 Y=\frac{y}{\sqrt{1-x}},\qquad
 u=2(x-t^2/2+iY),\qquad v=2(x-t^2/2-iY).
\end{equation}
The Jacobian in $(x,y)$ at the origin is invertible, and the equation
becomes exactly $uv=t^4$. The involution exchanges $u$ and $v$.

Perform the blowups first on the algebraic surface over $\mathbb A^1$
and then restrict to $\Delta$. In the $t$-chart of the first blowup,
$u=tu_1$, $v=tv_1$, and the strict transform is $u_1v_1=t^2$.
Its only singular point is its origin. In the $u$-chart, write
$v=uv_1$, $t=ut_1$. The equation is $v_1=u^2t_1^4$, a smooth chart
with coordinates $(u,t_1)$ and base function $t=ut_1$. The $v$-chart
is symmetric. Blow up the remaining origin. Its $t$-chart has
$u_1=ta$, $v_1=tb$ and $ab=1$. Its $u_1$-chart has
$v_1=u_1b$, $t=u_1a$ and $b=a^2$, with base function $t=u_1a$;
the other chart is symmetric. All charts are smooth and the zero divisor
of $t$ is reduced with ordinary crossings.

The four crossing charts can be written uniformly as
\begin{equation}\label{eq:resolution-charts}
 u=x_i^{i+1}y_i^i,\qquad
 v=x_i^{3-i}y_i^{4-i},\qquad t=x_iy_i,
 \qquad 0\le i\le3,
\end{equation}
with adjacent transition
$x_{i+1}=y_i^{-1}$, $y_{i+1}=x_iy_i^2$ wherever $y_i\ne0$.
They exhibit the chain of three exceptional rational curves. The original
central cubic is irreducible and nodal, with rational normalization. Its
proper transform joins the two ends of the chain at its two distinct
branches. This gives the asserted four-cycle. It is the usual resolution
pattern of the $A_3$ singularity; the charts additionally verify the
multiplicities of this particular base function.

Both singular-point centres are invariant under $\sigma$, so the universal
property of the blowup lifts the involution at each step. In
\eqref{eq:resolution-charts} it sends $(x_i,y_i)$ to $(y_i,x_i)$ in
chart $3-i$. Label the proper transform by $D_0$ and the exceptional
chain successively by $D_1,D_2,D_3$. The resulting component action is
$j\mapsto-j$. It fixes $D_0,D_2$ as sets, interchanging their incident
nodes, and exchanges $D_1,D_3$. Thus all four nodes are paired.

Blowups and the original plane-cubic morphism are projective and proper.
The resolved total space is smooth, and its nonconstant base coordinate
is a non-zero-divisor. It is therefore flat over the nonsingular curve
base. The blowups do not change any noncentral fibre. Shrinking the disc
if necessary gives the claimed family.
\end{proof}

Choose a smooth K\"ahler form on $\cE$, using a relative projective
embedding and a positive base form, and average it under $\sigma$.
Write the resulting invariant form as $\omega_{\cE}$ and put
\begin{equation}\label{eq:component-areas}
 A_j=\int_{D_j}\omega_{\cE}>0,\qquad A_3=A_1.
\end{equation}
No equality between $A_0,A_1,A_2$ is imposed. These actual areas will
enter the spectral mass and the finite-dimensional metric evolution.

\subsection{A holomorphic gluing family and its lift}
Choose disjoint node neighbourhoods $N_j$ at $D_j\cap D_{j+1}$ and
component neighbourhoods $O_j$ extending slightly into their incidence
collars. They can be chosen so that
$\sigma(O_j)=O_{-j}$ and $\sigma(N_j)=N_{-j-1}$. Each node set meets
its two component sets on separate connected collars, and there are no
triple overlaps. For example, choose two radii $0<a<b$ and restrict to
$|t|<a^2$: the two collars of a node are disjoint because $uv=t$.
Properness ensures that this cover of a neighbourhood of the compact
central fibre contains every fibre over a sufficiently small disc.

For $q\in\CC^*$ define
\begin{equation}\label{eq:r-j}
 R(q)=\begin{pmatrix}q&0\\0&q^{-1}\end{pmatrix},\qquad
 J=\begin{pmatrix}0&i\\i&0\end{pmatrix}=i\sigma_x.
\end{equation}
Here $\sigma_x$ is the first Pauli matrix. On $O_j\times\CC^*$ and
$N_j\times\CC^*$ take trivial rank-two bundles, with fibre coordinates
$w_j$ and $z_j$, respectively. Use transitions
\begin{equation}\label{eq:incidence-transitions}
 w_j=z_j,\qquad w_{j+1}=R(q)z_j
\end{equation}
on the two incidences.

\begin{proposition}\label{prop:flat-family}
The transitions \eqref{eq:incidence-transitions} define a holomorphic
rank-two bundle $V$ with trivialized determinant on $\cE\times\CC^*$.
For fixed $q$ it carries a flat holomorphic connection. On $|q|=1$ the
identity metrics glue to a parallel unitary metric. The maps
\begin{equation}\label{eq:projective-lift}
 w_{-j}=Jw_j,\qquad z_{-j-1}=R(q)^{-1}Jz_j
\end{equation}
give a lift of $\sigma$ whose square is $-I$. They preserve the connection,
and on the unit circle preserve its metric.
\end{proposition}

\begin{proof}
The transitions are holomorphic and invertible on the actual overlaps,
have determinant one, and have no triple-overlap condition to satisfy.
For each fixed $q$, they are constant in the $\cE$ coordinates, so the
connections $d$ in the patch frames agree. Flatness in the $q$ direction
is not asserted. Unitarity gives the metric statement.

The identities $JR=R^{-1}J$ and $J^2=-I$ verify the lift. The source
left incidence maps to the target right incidence, where
$Rz_{-j-1}=Jz_j=Jw_j=w_{-j}$. On the other incidence,
$Jw_{j+1}=JRz_j=R^{-1}Jz_j$, as required. Both component and node
lifts square to $-I$, since $(R^{-1}J)^2=-I$. The compatibility is
holomorphic in $q$ and preserves the flat connection for fixed $q$.
\end{proof}

For the unitary family write
\begin{equation}\label{eq:unitary-parameter}
 q=e^{i\theta/4},\qquad
 R=e^{i\theta\sigma_z/4},\qquad
 U_\theta=R^4=\begin{pmatrix}e^{i\theta}&0\\0&e^{-i\theta}\end{pmatrix},
 \quad \sigma_z=\begin{pmatrix}1&0\\0&-1\end{pmatrix}.
\end{equation}
A positively oriented path around the component cycle passes from the
$j$ frame to the $j+1$ frame by $R$ at every step. Its actual flat
holonomy is therefore $U_\theta$. The vanishing loop has trivial
holonomy, because the whole node patch carries the connection $d$.
These loops generate the elliptic fibre's fundamental group. Thus the
transition product has been identified with connection holonomy, with
an orientation and path specified.

\subsection{The free quotient and its metric}
Fix a projective K3 surface $S$ with a free Enriques involution $\iota$.
Such projective examples exist, and the involution acts as $-1$ on the
holomorphic two-form; see \cite[Example 2.10]{hashimoto-kanazawa}.
The smooth quotient construction below is the classical Enriques
Calabi--Yau example \cite[Example 3.5]{hashimoto-kanazawa}.
Choose an invariant K\"ahler class by averaging an ample class.
By \cite[Theorem 2]{yau}, its unique Ricci-flat representative is
$\iota$-invariant. Rescale it to obtain a form $\omega_S$ with
\begin{equation}\label{eq:k3-volume}
 \int_S\frac{\omega_S^2}{2}=1.
\end{equation}
The Chern connection on $T^{1,0}S$ has holonomy $\SU(2)$ by
\cite[Propositions 1 and 4]{beauville}. Its complex
parallel endomorphisms are consequently scalar by Schur's lemma. This
statement concerns complex-linear tangent endomorphisms, not the larger
real tangent commutant.
Fix the holomorphic two-form on $S$ to have pointwise norm one for this
metric. Such a scaling is possible because its norm is constant for a
Ricci-flat K\"ahler metric. Its dual therefore gives a unit-norm
trivialization of $\det T^{1,0}S$.

\begin{proposition}\label{prop:threefold-family}
The quotient
\begin{equation}\label{eq:threefold-family}
 \cX=(S\times\cE)/\langle\tau\rangle,
 \qquad \tau=(\iota,\sigma),
\end{equation}
is a proper projective flat family with smooth total space. Its central
fibre is a reduced simple normal-crossing threefold with three smooth
components and interval dual graph. For $t\ne0$, its smooth connected
fibre $X_t$ has trivial canonical bundle, $h^{1,0}=h^{2,0}=0$, and an
index-two subgroup $\ZZ^2$ in its fundamental group. The form
$\omega_S+\omega_{\cE}|_{E_t}$ descends to a K\"ahler form $\omega_t$ on
$X_t$, where $E_t$ is the elliptic fibre.
\end{proposition}

\begin{proof}
The involution is free on the entire total space because $\iota$ is free.
The quotient map is finite \emph{\'{e}tale}; hence the quotient is smooth
and the local crossing equation is unchanged. The fixed component sets
give $(S\times D_0)/\langle\tau\rangle$ and
$(S\times D_2)/\langle\tau\rangle$. The exchanged pair
$S\times D_1,S\times D_3$ has image isomorphic to $S\times\PP^1$.
The two pairs of nodes give two smooth double loci isomorphic to $S$.
No node is fixed and no branch of a node is identified with the other
branch. These observations prove the asserted simple normal crossings
and interval incidence.

For projectivity choose a relatively ample line bundle $L$ on
$S\times\cE$, using an ample line on $S$ and a relative embedding of
$\cE$. The line $L\otimes\tau^*L$ has the order-two linearization which
exchanges its factors. It descends under the free action; its descent
is relatively ample because ampleness can be tested after finite
surjective pullback. Properness follows from the finite quotient of the
proper family. Flatness follows from the non-zero-divisor $t$ on the
smooth total space.

On each smooth fibre both the K3 two-form and the elliptic one-form
change sign under their respective involutions. Their product descends
to a nowhere vanishing three-form. K\"unneth and invariants under the
finite \emph{\'{e}tale} cover give $h^{1,0}=h^{2,0}=0$.
The cover $S\times E_t$ has fundamental group $\ZZ^2$, which injects
with index two into the quotient fundamental group. The product K\"ahler
form is invariant and hence descends.
\end{proof}

The elliptic metric is the restriction of a smooth total-space metric.
It is generally not flat, so the descended threefold metric is generally
not Ricci-flat. All subsequent analytic statements use this particular
metric, including its actual component areas \eqref{eq:component-areas}.

\subsection{Rank-four descent, determinant and second Chern class}
On $S\times\cE\times\CC^*$ form
\begin{equation}\label{eq:cover-bundle}
 \widetilde{\cB}=p_S^*T^{1,0}S\otimes p_{\cE}^*V.
\end{equation}
The tangent map $i\,d\iota$ has square $-I$. Tensor it with the lift
\eqref{eq:projective-lift} of $\sigma$.

\begin{proposition}\label{prop:bundle-descent}
The tensor lift defines an order-two equivariant structure on
$\widetilde{\cB}$ and gives a holomorphic rank-four bundle $\cB$ on
$\cX\times\CC^*$. Its determinant has an actual holomorphic
trivialization. For every $t\ne0$ and every $q$, a specified embedded
K3 slice $S_e\subset X_t$ satisfies
\begin{equation}\label{eq:chern-pairing}
 \int_{S_e}c_2(\cB_{t,q}|_{S_e})=48.
\end{equation}
\end{proposition}

\begin{proof}
The two factors of the lift square to $-I$, so their tensor product
squares to $I$. Incidence compatibility follows from
Proposition~\ref{prop:flat-family}. Thus the tensor bundle has genuine
descent under the free involution, holomorphically in $q$.

Trivialize $\det T^{1,0}S$ by the dual of its holomorphic two-form.
Anti-symplecticity gives $\det(d\iota)=-1$ in these trivializations,
and hence $\det(i\,d\iota)=1$. The matrices $J$ and $R^{-1}J$ have
determinant one. Since both tensor factors have rank two,
\[
 \det(A\otimes B)=(\det A)^2(\det B)^2=1.
\]
The determinant trivialization is invariant under the descent datum
and therefore descends; this is stronger than a torsion determinant line.

Choose $e\in E_t$ which is not a two-torsion point. The image of
$S\times\{e\}$ is an embedded K3 slice: no two of its points are
$\tau$-related because $-e\ne e$, and the proper map is locally an
embedding under the \emph{\'{e}tale} quotient. Restriction of the bundle
to that slice is $T^{1,0}S\oplus T^{1,0}S$. The identities
$c_1(T^{1,0}S)=0$ and $\int_Sc_2(T^{1,0}S)=24$ give
\eqref{eq:chern-pairing}. No division of an integral class by the cover
degree is used.
\end{proof}

\subsection{The Hermitian--Einstein connection and stability}
Let $h_S$ be the tangent metric induced by $\omega_S$. On the unitary
locus equip \eqref{eq:cover-bundle} with $h_S\otimes I$.
The lift is unitary, so this metric and its Chern connection descend.
Write their descents as $h_0$ and $D_\theta$.

\begin{proposition}\label{prop:he-stability}
For the metric $\omega_t$ of Proposition~\ref{prop:threefold-family},
the reference metric $h_0$ is determinant-normalized Hermitian--Einstein
with zero Einstein constant. The bundle $\cB_{t,\theta}$ is slope stable
for $\theta\notin\pi\ZZ$ and strictly polystable for $\theta\in\pi\ZZ$.
Its complex parallel endomorphism algebra is
\begin{equation}\label{eq:parallel-algebra}
 \End_{\rm par}(\cB_{t,\theta})
 \simeq\Comm(U_\theta)\cap\Comm(J).
\end{equation}
\end{proposition}

\begin{proof}
The chosen determinant trivialization has unit norm for $h_S$, and the
flat multiplicity determinant has unit norm for $I$. Hence their tensor
determinant, and its invariant descent, have unit norm for $h_0$.
The product connection has curvature $F_{T^{1,0}S}\otimes I$, since
the elliptic factor is flat. Product contraction gives
\[
 i\Lambda_{\omega_S+\omega_{\cE}}F_{\widetilde{\cB}}
 =(i\Lambda_{\omega_S}F_{T^{1,0}S})\otimes I=0.
\]
This uses Ricci flatness of the K3 metric and imposes no flatness condition
on the elliptic metric. It descends to the stated connection on $X_t$.

The K3 holonomy makes every parallel complex endomorphism on the cover
equal to $I_{T^{1,0}S}\otimes M$ for $M\in\Comm(U_\theta)$.
Compatibility with the lift at the $D_0$ component is exactly
$MJ=JM$. This proves \eqref{eq:parallel-algebra}; the other component
and node conditions follow by transport and the verified lift.

For $\theta\notin\pi\ZZ$, $U_\theta$ has distinct eigenvalues.
Commutation with it makes $M$ diagonal, and commutation with $J$ makes
its diagonal entries equal. The parallel algebra is scalar. The
Hermitian--Einstein polystable decomposition
\cite[Theorem 5.8.3]{kobayashi} therefore has only one
stable summand, because any further summand would give a nonscalar
parallel projection. Thus the bundle is stable. At $\theta\in\pi\ZZ$
the matrix $U_\theta$ is central and $\Comm(J)\simeq\CC\oplus\CC$.
Its eigenprojections split the bundle into two rank-two
Hermitian--Einstein summands, proving strict polystability. Their slopes
are equal and the total slope is zero for the actual $\omega_t$.
\end{proof}

The holomorphic family is defined for $q\in\CC^*$, whereas the displayed
Hermitian--Einstein connections and stability analysis concern $|q|=1$.
The approach to strict polystability is a genuine change of holomorphic
bundle gluing, independent of $t$. No nonunitary or unstable locus is
asserted by this construction.
